\section{A sharp pairwise inverse through a triple collision}
\label{sec:v5-pairwise}
We now compare two arbitrary nearby curvature triples, rather than a
triple with the circular reference alone. In the ambient model with an independent area normalizer, four
amplitudes give an inverse without separating repeated nodes. The
constrained physical family admits a three-amplitude refinement near
$R=1/4$, proved in Section~\ref{sec:v6-three}.

Fix $g,R,A_0>0$ and consider the observation model
\begin{equation}\label{eq:v5-amplitude-model}
 C_j(A,r)=\frac1A\sum_{i=1}^3\Phi_j(r_i),\qquad
 \Phi_j(r)=\csch\bigl(j\arcosh(1+g/r)\bigr).
\end{equation}
Here $A$ is near $A_0$ and $r_i$ are near $R$. The equal-gap geometric
family of Theorem~\ref{thm:g-inverse} is a subfamily of this model, with
$r_i=\kappa_i^{-1}$ and its area given by \eqref{eq:v3-area}. Allowing
$A$ to vary independently in the observation model proves a stronger
nuisance-parameter statement; it does not make area an independent
shape coordinate in that geometric family. When $r_i^{-1}$ is
interpreted as an actual contact curvature, each channel is assumed to
have equal facing curvatures. Without that equality, the same scalar
model describes the effective channel curvatures
\eqref{eq:v6-effective-curvature}, not both endpoint curvatures.

\begin{theorem}[Four-amplitude pairwise stability]\label{thm:v5-pairwise}
There is a neighborhood of $(A_0,R,R,R)$ on which, for every two
members of \eqref{eq:v5-amplitude-model},
\begin{equation}\label{eq:v5-pairwise-bound}
 |A-\widetilde A|\le C\delta_4,\qquad
 \operatorname{dist}_{\rm match}(\kappa,\widetilde\kappa)
                         \le C\delta_4^{1/3},\qquad
 \delta_4=\max_{1\le j\le4}|C_j-\widetilde C_j|.
\end{equation}
The inverse determines multiplicities and requires neither channel
labels nor a lower bound on the distance between distinct curvatures.
In particular the first four amplitudes determine the area and unordered
triple locally, even at complete coalescence. With deterministic error
at most $\epsilon$ in these four amplitudes, a reconstruction constrained
to a sufficiently small compact neighborhood has area error $C\epsilon$
and matching curvature error $C\epsilon^{1/3}$.
\end{theorem}

\begin{proof}
Put $x_i=r_i-R$ and let $e_1,e_2,e_3$ be their elementary symmetric
polynomials. Define
\[
 p_e(z)=z^3-e_1z^2+e_2z-e_3.
\]
For a small fixed circle $|z|=b$ around zero, $\Phi_j(R+z)$ is analytic
inside it, with the positive-real branch at $R$. If $e$ is sufficiently
small, all roots of $p_e$ lie strictly inside that circle. The symmetric
sum in \eqref{eq:v5-amplitude-model} has the analytic extension
\begin{equation}\label{eq:v5-contour}
 M_j(e)=\frac1{2\pi i}\int_{|z|=b}
             \Phi_j(R+z)\frac{p_e'(z)}{p_e(z)}\dd z.
\end{equation}
This formula counts roots with multiplicity and defines an analytic
function on a full coefficient neighborhood, not merely on its
real-rooted part. It is real for real $e$. Expanding the Newton sums at
$e=0$ gives
\begin{equation}\label{eq:v5-symmetric-jets}
 M_j(0)=3\Phi_j(R),\quad
 \partial_{e_1}M_j=\Phi_j'(R),\quad
 \partial_{e_2}M_j=-\Phi_j''(R),\quad
 \partial_{e_3}M_j=\tfrac12\Phi_j'''(R)
 \quad(e=0).
\end{equation}
For clarity, the first three Newton sums are
$e_1$, $e_1^2-2e_2$, and $e_1^3-3e_1e_2+3e_3$; all sums of degree
at least four have zero first differential at $e=0$. This also follows
by differentiating the contour integral.

Consider the map
$\Psi:(A,e_1,e_2,e_3)\mapsto(A^{-1}M_j(e))_{j=1}^4$.
Its Jacobian determinant at $(A_0,0,0,0)$ is
\begin{equation}\label{eq:v5-jacobian}
 \det D\Psi=\frac3{2A_0^5}\,
   \det\bigl[\Phi_j^{(k)}(R)\bigr]_{1\le j\le4,\,0\le k\le3}>0.
\end{equation}
We verify the strict sign explicitly. Put $c=1+g/R>1$ and write
$\Phi_j(r)=w(c(r))f_j(c(r))$, where $c(r)=1+g/r$,
$w(c)=(c^2-1)^{-1/2}$, and
\[
 f_1=1,\qquad f_2=\frac1{2c},\qquad
 f_3=\frac1{4c^2-1},\qquad f_4=\frac1{4c(2c^2-1)}.
\]
Column elimination in the product and chain rules gives
\[
 \det[\Phi_j^{(k)}(R)]
  =\left(\frac g{R^2}\right)^6(c^2-1)^{-2}
        \det[f_j^{(k)}(c)].
\]
The last determinant reduces, using $f_1=1$, to the three by three
matrix of the first, second, and third derivatives of $f_2,f_3,f_4$.
Differentiating these displayed rational functions and putting the
entries over their common denominator yields
\begin{equation}\label{eq:v5-wronskian}
 \det[f_j^{(k)}(c)]
 =\frac{12(64c^8+32c^6+116c^4+4c^2+1)}
 {c^4(4c^2-1)^4(2c^2-1)^4}.
\end{equation}
Every factor is strictly positive for $c>1$. The column factors in
\eqref{eq:v5-symmetric-jets} are $-3/A_0^2$, $1/A_0$, $-1/A_0$ and $1/(2A_0)$,
which proves the factor $3/(2A_0^5)$ in \eqref{eq:v5-jacobian}.
Thus the analytic inverse-function theorem applies in the full
four-dimensional coefficient space. Choose a convex ball in its data
image and then a smaller parameter neighborhood whose image lies in
that ball. A bounded derivative of the inverse on the ball gives
\begin{equation}\label{eq:v5-coefficient-stability}
 |A-\widetilde A|+\sum_{k=1}^3|e_k-\widetilde e_k|
                           \le C\delta_4
\end{equation}
for any two members of this smaller neighborhood, including its
real-rooted boundary strata. This is the step for which using the
symmetric coordinates, rather than differentiating individually
labelled roots, is essential.

We give the needed root estimate with multiplicities. For two monic
cubics $p,q$ whose roots lie in a fixed bounded set and whose coefficient
distance is $\epsilon$, one has $|p(z)-q(z)|\le C_0\epsilon$ on a common
bounded disk containing the small neighborhoods below. Set
$\rho=(2C_0\epsilon)^{1/3}$. On the boundary of the union of the open
radius-$\rho$ disks about the three roots of $p$, every root distance
is at least $\rho$. Hence $|p(z)|\ge\rho^3>|p(z)-q(z)|$ there.
Rouch\'e's theorem on each connected component shows that the two
polynomials have the same number of roots in that component, counted
with multiplicities. If two boundary arcs are tangent, perturb the
radius slightly and pass to the limit; the strict inequality persists.
Every component consists of at most three intersecting disks and has
diameter at most $6\rho$. Match the roots inside each component. This
proves a matching distance at most $C\epsilon^{1/3}$. At $\epsilon=0$
the polynomials coincide. The assertion for bounded, not necessarily
small, $\epsilon$ follows by enlarging the constant.

Apply this estimate to \eqref{eq:v5-coefficient-stability}. The radii
stay in a compact interval bounded away from zero, so $r\mapsto r^{-1}$
is Lipschitz there. This proves \eqref{eq:v5-pairwise-bound}. For noisy
data, minimize the maximum discrepancy over a small compact admissible
parameter set containing the true member. A minimum exists by
continuity. Its exact amplitude vector differs from the true one by at
most $2\epsilon$, since the true member is itself feasible. Apply the
pairwise estimate. Uniqueness of this minimizer is not required for the
error bound, nor is an efficient numerical optimization algorithm being
assumed in its proof.
\end{proof}

\subsection{Sharpness in the physical geometric family}
For $a\ge0$ use the same absolute-amplitude topology as before,
$\|D\|_a=\sup_{j\ge1}e^{aj}|D_j|$. Choose a small fixed-$R$
neighborhood in \eqref{eq:g-isogap-family} on which all exponents exceed
some $\gamma_->a$.

\begin{theorem}[Optimal pairwise exponent]\label{thm:v5-sharpness}
On a sufficiently small neighborhood in that physical family,
\[
 \operatorname{dist}_{\rm match}(\kappa,\widetilde\kappa)
       \le C_a\|C-\widetilde C\|_a^{1/3}.
\]
The exponent $1/3$ cannot be increased, even if the entire sequence is
observed. On the two-sided path $\alpha=\zeta=0$, $\beta=s$,
\begin{equation}\label{eq:v5-cubic-path}
 \|C(s)-C(-s)\|_a\asymp |s|^3,\qquad
 \operatorname{dist}_{\rm match}(\kappa(s),\kappa(-s))\asymp |s|.
\end{equation}
The circular-reference exponent $1/2$ of
Theorem~\ref{thm:v4-sequence-stability} remains valid and sharp for
its different, reference-point quantifiers.
\end{theorem}
\begin{proof}
The first assertion follows from Theorem~\ref{thm:v5-pairwise}, since
the first four errors are bounded by the sequence norm. For sharpness
use exactly the realized path of Proposition~\ref{prop:v3-area-coalescence}:
\[
 r(s)=(R+36s,R-18s,R-18s),\qquad A(s)=A_0+45\pi s^2/4.
\]
Here $A_0$ is the actual circular free area. Differentiation of
$2e^{-j\gamma}/(1-e^{-2j\gamma})$, together with the bounded first three
derivatives of $\gamma(r)$ and $A(s)^{-1}$, gives
\[
 \sup_{|s|\le s_*}|C_j'''(s)|
                     \le B(1+j)^3e^{-j\gamma_-}.
\]
The first derivative at zero vanishes. Subtract the Taylor expansions
at $s$ and $-s$: their constant and quadratic terms cancel, as does the
zero linear term. The strict margin $a<\gamma_-$ absorbs the polynomial
in $j$, proving $\|C(s)-C(-s)\|_a\le B_a|s|^3$.

The first amplitude supplies a nonzero lower bound. Direct differentiation
gives
\[
 \Phi_1(r)=\frac r{\sqrt{g(g+2r)}},\qquad
 \Phi_1'''(r)=\frac{3(3g+r)}{\sqrt g\,(g+2r)^{7/2}}>0.
\]
Since $36^3+2(-18)^3=34992$ and $A(s)$ is even,
\begin{equation}\label{eq:v5-cubic-coefficient}
 C_1(s)-C_1(-s)
       =\frac{11664\Phi_1'''(R)}{A_0}s^3+O(s^5).
\end{equation}
This proves the other sequence bound. In one real dimension sorted
matching minimizes the maximum coordinate distance: an inverted pair
can be uncrossed without increasing either extreme distance. Sorting
the two curvature triples, their middle pair is the largest difference
for all sufficiently small $s>0$, and gives
\[
 \operatorname{dist}_{\rm match}(\kappa(s),\kappa(-s))
             =\frac{36s}{R^2-324s^2}.
\]
The negative case follows by symmetry. If a pairwise exponent
$\theta>1/3$ held uniformly, these estimates would imply
$b|s|\le C|s|^{3\theta}$ for all small nonzero $s$, a contradiction.
This proves sharpness without replacing the infinite sequence by a
finite numerical prefix.
\end{proof}

The reference and pairwise results describe different singularities of
one observation map. Distance from the circular point first detects
quadratic dispersion. Comparing two noncircular triples on opposite
sides of it can cancel that dispersion as well, leaving a cubic term.
The four-amplitude inverse proves that the latter cancellation is the
worst local pairwise scale in the stated three-channel model. This
positive conclusion does not use relative accuracy at arbitrarily small
amplitudes or a separation-dependent Prony bound.
